\section{本章小结}
到这里，我们把“接收到了无线电数据”拆成了一条能逐步检查的链：硬件留下 IQ，数字处理找到参考并恢复比特，长度、CRC 和命令确认候选帧，字段再按单位、坐标系和时间进入上层。离线项目让你固定输入、修改参数、重放坏帧；它不能替代天线、调谐和真实波形的检查。先练熟录包、同步和拒收这几步，再按设备与协议文档补齐实际接收机。

\begin{problemset}
\renewcommand{\theenumi}{\arabic{enumi}}
\item\label{prob:v2-sdr-1} 基础：采样前已被模拟滤波器去掉的频段，能否在 IQ 回放时恢复？ADC 饱和又能否只换软件补好？
\item\label{prob:v2-sdr-2} 基础：IQ 为 $3+4j$，幅值是多少？若只保存幅值，正负频率与相位还剩下什么信息？
\item\label{prob:v2-sdr-3} 基础：采样率 8000 Hz、每符号 8 个样本，BPSK 的符号率与比特率分别是多少？本例一帧含前导但不含空样本，要几个样本？
\item\label{prob:v2-sdr-4} 提高：哈希一致但元数据的采样率写错，为什么回放仍可能错？给出比文件哈希更完整的录包交付约定。
\item\label{prob:v2-sdr-5} 基础：同步、符号判决和 CRC 分别解决什么？CRC 成功能否证明字段属于正确设备、坐标系和当前时刻？
\item\label{prob:v2-sdr-6} 提高：未知频偏与恒定相位有什么区别？为什么前导估计出的一个相角不能一直抵消频偏？
\item\label{prob:v2-sdr-7} 提高：切换到 information 后收到 level2 的迟到解码结果，应该检查什么？仅将状态名改过来是否已经完成实机切换？
\item\label{prob:v2-sdr-8} 基础：同一时钟下 now=20.0 s、radio=19.6 s、vision=19.9 s、有效期 0.25 s，选择哪一路？两路都过期或时间戳在未来时怎样处理？
\item\label{prob:v2-sdr-9} 实践：用 \texttt{--phase -1.2 --offset-hz -125 --x-mm 3000} 在新目录生成并回放；记录起点、字段、采样元数据，再故意损坏帧并展示拒收。写清它验证了哪些假设，尚未验证什么。
\end{problemset}

\section{本章验收}
\begin{enumerate}
\item 在两次独立进程中生成、保存、读取和解码录包；检查者更改相位、毫米坐标或合法频偏，能指出输入与结果的对应关系。
\item 对正常、CRC 损坏和截断样本分别演示处理结果，指出同步与帧校验的边界；不把理想单帧实验称为实战接收机。
\item 按检查者给出的事件序列展示配置阶段和代次，再模拟无线电过期、视觉可用及两路不可用，解释输出状态。
\item 记录已检查的硬件依赖与未验证项；有获准设备时现场接收一个缓冲、检查元数据并清理，说明真实波形解码所需的协议与同步条件。
\end{enumerate}
